% Report-ready material generated from experiment 20260908_botnet_excluded_v1.
% Source of truth: data_lineage.json, model_selection.json, and final_evaluation.json.

% --- Replacement text for Section 3.1.2 ---
The experimental corpus was generated by re-extracting the CICIDS2017
PCAP files with NFStream. The labeled source contained 3,351,784 flows
and 57 numerical flow features. Botnet (735 flows), Infiltration (72
flows), and Heartbleed (one flow) were excluded from supervised machine
learning because their post-extraction support was insufficient for
reliable stratified validation and class-wise generalization analysis.
No invalid numerical row remained after the scope filter. Deduplication
reduced the BENIGN class from 2,889,264 to 1,188,807 distinct feature
vectors, removing 1,700,457 duplicates. The retained observations were
mapped to five classes: BENIGN, Brute Force, DoS/DDoS, PortScan, and Web
Attack.

To prevent the previously inspected random hold-out partition from being
reused as a final test, a chronological 80/20 split was performed
independently within each mapped class. Observations were ordered by
their NFStream first-seen timestamp; the earliest observations formed the
development partition and the latest observations formed the final test
partition. All observations sharing a boundary timestamp were assigned
to the final test side. The resulting development and final-test sets
contained 1,320,394 and 330,125 flows, respectively. Timestamp and
network identifiers were used only to construct the split and were not
provided to the models.

\begin{table}[H]
  \caption{Machine Learning Dataset Class Distribution}
  \centering
  \small
  \begin{tabularx}{\textwidth}{|X|r|r|r|r|}
    \hline
    \textbf{Class} & \textbf{Development} & \textbf{Ratio} & \textbf{Final test} & \textbf{Ratio} \\ \hline
    BENIGN & 951,045 & 72.03\% & 237,762 & 72.02\% \\ \hline
    Brute Force & 5,106 & 0.39\% & 1,277 & 0.39\% \\ \hline
    DoS/DDoS & 234,920 & 17.79\% & 58,735 & 17.79\% \\ \hline
    PortScan & 127,627 & 9.67\% & 31,926 & 9.67\% \\ \hline
    Web Attack & 1,696 & 0.13\% & 425 & 0.13\% \\ \hline
    \textbf{TOTAL} & \textbf{1,320,394} & \textbf{100.00\%} & \textbf{330,125} & \textbf{100.00\%} \\ \hline
  \end{tabularx}
\end{table}

% --- Replacement text and table for Section 3.2.1 ---
Model selection was performed exclusively on the development partition.
For every Optuna trial, feature importance was fitted on each training
fold, the minimum feature subset reaching 95\% cumulative importance was
selected, and SMOTE was applied only to Brute Force and Web Attack
observations in that training fold. The untouched validation fold was
then scored using macro-averaged F1. A seeded TPE sampler
(\texttt{seed=42}) completed 30 trials for each algorithm using the same
stratified five-fold partitions. The model family was selected by the
highest mean cross-validation macro-F1; the final test was not used for
algorithm, hyperparameter, or feature selection.

\begin{table}[H]
  \caption{Hyperparameter Search Spaces and Cross-Validation Performance}
  \centering
  \scriptsize
  \begin{tabularx}{\textwidth}{|p{1.8cm}|X|X|p{1.8cm}|p{1.8cm}|}
    \hline
    \textbf{Algorithm} & \textbf{Search space} & \textbf{Best configuration} & \textbf{Features per fold} & \textbf{CV macro-F1} \\ \hline
    LightGBM & max depth 3--15; leaves 8--128; learning rate 0.01--0.2; estimators 50--150; subsample and column sample 0.6--1.0 & depth 14; leaves 53; learning rate 0.055585; estimators 117; subsample 0.820544; column sample 0.607129 & 5, 5, 5, 5, 5 & $0.964517 \pm 0.003156$ \\ \hline
    Random Forest & max depth 5--25; estimators 50--150; minimum split 2--10; minimum leaf 1--5 & depth 25; estimators 90; minimum split 6; minimum leaf 1 & 36, 36, 35, 34, 35 & $0.987258 \pm 0.001313$ \\ \hline
    XGBoost & max depth 3--10; learning rate 0.01--0.2; estimators 50--150; subsample and column sample 0.6--1.0 & depth 8; learning rate 0.154810; estimators 78; subsample 0.854758; column sample 0.696565 & 20, 19, 17, 18, 18 & $\mathbf{0.994152 \pm 0.000673}$ \\ \hline
  \end{tabularx}
\end{table}

XGBoost was selected before final-test evaluation because it achieved the
highest development cross-validation macro-F1. After model-family and
hyperparameter selection, feature importance was refitted using the full
development partition. The 95\% cumulative-importance rule selected 21
features for the locked final models.

\begin{table}[H]
  \caption{Twenty-One-Feature Vector Selected for the Final XGBoost Models}
  \centering
  \small
  \begin{tabularx}{\textwidth}{|r|X|r|X|}
    \hline
    1 & Bwd Packets/s & 12 & Flow Bytes/s \\ \hline
    2 & Max Packet Length & 13 & Total Length of Bwd Packets \\ \hline
    3 & RST Flag Count & 14 & Fwd IAT Min \\ \hline
    4 & Bwd Packet Length Std & 15 & Packet Length Mean \\ \hline
    5 & Total Fwd Packets & 16 & Fwd PSH Flags \\ \hline
    6 & Down/Up Ratio & 17 & Total Backward Packets \\ \hline
    7 & SYN Flag Count & 18 & Packet Length Std \\ \hline
    8 & PSH Flag Count & 19 & Fwd Packet Length Max \\ \hline
    9 & Total Length of Fwd Packets & 20 & Init\_Win\_bytes\_forward \\ \hline
    10 & ACK Flag Count & 21 & Fwd Packet Length Std \\ \hline
    11 & Flow IAT Std & & \\ \hline
  \end{tabularx}
\end{table}

% --- Replacement text and tables for Section 4.1.1 ---
After XGBoost, its hyperparameters, and the 21-feature subset had been
locked, the final test partition was used to generate predictions. The
binary component was evaluated at the classifier's default decision
boundary. It achieved 99.6483\% accuracy and a weighted F1-score of
99.6489\%, with an FPR of 0.4492\% and an FNR of 0.1007\%.

\begin{table}[H]
  \caption{Binary Active Model Confusion Matrix}
  \centering
  \small
  \begin{tabularx}{\textwidth}{|X|r|r|r|}
    \hline
    \textbf{Actual / Predicted} & \textbf{Predicted benign} & \textbf{Predicted malicious} & \textbf{Total} \\ \hline
    Actual benign & 236,694 & 1,068 & 237,762 \\ \hline
    Actual malicious & 93 & 92,270 & 92,363 \\ \hline
    \textbf{Total} & \textbf{236,787} & \textbf{93,338} & \textbf{330,125} \\ \hline
  \end{tabularx}
\end{table}

The five-class component achieved 99.6398\% accuracy, a macro-F1 of
99.0701\%, and a weighted F1-score of 99.6408\% on the same 330,125
final-test flows.

\begin{table}[H]
  \caption{Multiclass Classification Performance Report}
  \centering
  \small
  \begin{tabularx}{\textwidth}{|X|r|r|r|r|}
    \hline
    \textbf{Class} & \textbf{Precision} & \textbf{Recall} & \textbf{F1-score} & \textbf{Support} \\ \hline
    BENIGN & 0.99952 & 0.99554 & 0.99753 & 237,762 \\ \hline
    Brute Force & 1.00000 & 0.97807 & 0.98892 & 1,277 \\ \hline
    DoS/DDoS & 0.98224 & 1.00000 & 0.99104 & 58,735 \\ \hline
    PortScan & 0.99981 & 0.99721 & 0.99851 & 31,926 \\ \hline
    Web Attack & 0.98333 & 0.97176 & 0.97751 & 425 \\ \hline
    \textbf{Macro average} & \textbf{0.99298} & \textbf{0.98852} & \textbf{0.99070} & \textbf{330,125} \\ \hline
    \textbf{Weighted average} & \textbf{0.99645} & \textbf{0.99640} & \textbf{0.99641} & \textbf{330,125} \\ \hline
  \end{tabularx}
\end{table}

The deployed ML decision rule evaluates both XGBoost models for every
flow and defines the attack score as the maximum of the binary attack
probability and one minus the multiclass BENIGN probability. At the
configured threshold of 0.85, this rule produced 236,704 true negatives,
1,058 false positives, 333 false negatives, and 92,030 true positives.
Its attack FPR was 0.4450\%, its attack FNR was 0.3605\%, and its
end-to-end five-class macro-F1 was 97.1550\%. The higher FNR relative to
the standalone binary classifier was principally associated with Brute
Force recall decreasing from 0.9781 to 0.8097. Therefore, threshold 0.85
should be presented as the current operational trade-off rather than as
an independently optimized threshold.

Bulk inference was benchmarked after one warm-up using five repetitions
over all 330,125 final-test rows. The median amortized prediction times
were 0.000189 ms/flow for the binary component and 0.000784 ms/flow for
the multiclass component. These measurements exclude NFStream flow
completion, feature extraction, serialization, Suricata correlation,
and firewall execution, and therefore must not be described as
end-to-end attack-detection latency.

% --- Reproducibility material suitable for an appendix ---
\begin{verbatim}
python -u src\run_botnet_excluded_experiment.py --stage prepare \
  --experiment-id 20260908_botnet_excluded_v1
python -u src\run_botnet_excluded_experiment.py --stage tune \
  --experiment-id 20260908_botnet_excluded_v1 --trials 30 --rf-jobs -1
python -u src\run_botnet_excluded_experiment.py --stage train-evaluate \
  --experiment-id 20260908_botnet_excluded_v1
python -u src\run_botnet_excluded_experiment.py --stage verify \
  --experiment-id 20260908_botnet_excluded_v1
\end{verbatim}

The input SHA-256 was
\texttt{fd277e4a6a8dd8088b724f80d0fca6d8c3d7f7453d4ad525e8eff10b85a5f579}.
The development and final-test SHA-256 values were
\texttt{aa69beb5498998386c6b86b1228f74517a4565d9f2a57513ee795617635cb2f7}
and
\texttt{ef6ce34df176b5af02d7df3e985c0df1ce3e6664e46a61a5e92ea232af4ceac6},
respectively. Exact package versions, commands, trial histories, fold
scores, model metadata, and artifact hashes are retained under experiment
identifier \texttt{20260908\_botnet\_excluded\_v1}.
